\uuid{j6aO}
\exo7id{4323}
\titre{exo7 4323}
\auteur{quercia}
\organisation{exo7}
\datecreate{2010-03-12}
\isIndication{false}
\isCorrection{false}
\chapitre{Intégration}
\sousChapitre{Intégrale de Riemann dépendant d'un paramètre}
\cpge{mpsi-mp2i.integration.c, pcsi.integration.b, ptsi.integration.b, mpsi-mp2i.nombres-reels-suites.d, pcsi.nombres-reels-suites.c, ptsi.nombres-reels-suites.c, bcpst1.integration-segment.a, bcpst1.integration-segment.b, bcpst1.suites-reelles}
\module{Analyse}
\niveau{L2}
\difficulte{}

\contenu{
\question{
Soit $I_n =  \int_{t=0}^1 t^n\ln(1+t^2)\,d t$.
Montrer que $I_n \to 0$ lorsque $n\to\infty$.
}
}
